% Standalone problem document, generated from the adjacent README.md.
% Compile with XeLaTeX. Mathematical content is maintained in Markdown.
\documentclass[11pt,a4paper]{article}
\usepackage[fontsize=10.5pt]{scrextend}
\usepackage[margin=22mm,headheight=15pt,headsep=9mm,footskip=11mm]{geometry}
\usepackage{fontspec}
\setmainfont{texgyrepagella}[Extension=.otf,UprightFont=*-regular,BoldFont=*-bold,ItalicFont=*-italic,BoldItalicFont=*-bolditalic]
\setsansfont{texgyreheros}[Extension=.otf,UprightFont=*-regular,BoldFont=*-bold,ItalicFont=*-italic,BoldItalicFont=*-bolditalic]
\setmonofont{texgyrecursor}[Extension=.otf,UprightFont=*-regular,BoldFont=*-bold,ItalicFont=*-italic,BoldItalicFont=*-bolditalic,Scale=MatchLowercase]
\usepackage{amsmath,amssymb}
\usepackage{unicode-math}
\setmathfont{texgyrepagella-math.otf}
\usepackage{xcolor,microtype,enumitem,needspace,fancyhdr}
\usepackage{bookmark}
\definecolor{ink}{HTML}{17344B}
\definecolor{accent}{HTML}{197D80}
\definecolor{muted}{HTML}{596773}
\hypersetup{colorlinks=true,linkcolor=accent,urlcolor=accent,pdftitle={IE-28 - Positive
diagonal preconditioning with a nilpotent stiff
limit},pdfauthor={Open Problems in Numerical Linear Algebra}}
\urlstyle{same}
\setlength{\parindent}{0pt}
\setlength{\parskip}{5pt plus 1pt minus 1pt}
\linespread{1.04}
\setlength{\emergencystretch}{3em}
\widowpenalty=10000
\clubpenalty=10000
\displaywidowpenalty=10000
\allowdisplaybreaks[1]
\setlist{nosep,leftmargin=1.5em}
\providecommand{\tightlist}{\setlength{\itemsep}{0pt}\setlength{\parskip}{0pt}}
\providecommand{\pandocbounded}[1]{#1}
\pagestyle{fancy}
\fancyhf{}
\fancyhead[L]{\sffamily\footnotesize\color{muted}OPEN PROBLEMS IN NUMERICAL LINEAR ALGEBRA}
\fancyhead[R]{\sffamily\bfseries\footnotesize\color{accent}IE-28}
\fancyfoot[L]{\parbox[b]{0.9\textwidth}{\sffamily\fontsize{8}{10}\selectfont\color{muted}Editorial ratings; a bounded search cannot certify that a problem remains unsolved.\\Literature check: 2026-09-11}}
\fancyfoot[R]{\sffamily\footnotesize\color{muted}\thepage}
\renewcommand{\headrulewidth}{0.3pt}
\renewcommand{\headrule}{\hbox to\headwidth{\color{accent}\leaders\hrule height \headrulewidth\hfill}}
\makeatletter
\renewcommand{\section}{\Needspace{5\baselineskip}\@startsection{section}{1}{0pt}{12pt plus 2pt minus 2pt}{4pt}{\sffamily\bfseries\large\color{ink}}}
\renewcommand{\subsection}{\Needspace{4\baselineskip}\@startsection{subsection}{2}{0pt}{9pt plus 1pt minus 1pt}{3pt}{\sffamily\bfseries\normalsize\color{ink}}}
\makeatother
\setcounter{secnumdepth}{0}
\begin{document}
{\sffamily\small\color{muted}Linear systems and elimination\par}
\vspace{3pt}
{\raggedright\hyphenpenalty=10000\exhyphenpenalty=10000\sffamily\bfseries\fontsize{21}{24}\selectfont\color{ink}Positive
diagonal preconditioning with a nilpotent stiff limit\par}
\vspace{7pt}
{\sffamily\small\textbf{Difficulty:} challenging\quad\textbf{Importance:} interesting
to specialist\par}
{\sffamily\small\bfseries\color{accent}Status: Partially resolved\par}
\vspace{4pt}
{\color{accent}\hrule height 0.8pt}
\vspace{6pt}
\textbf{Topic:} Diagonal preconditioners for collocation systems

\textbf{Rating rationale:} The problem is a uniform existence assertion
for a positive diagonal solution of nonlinear spectral equations. It
underlies parallel iteration for implicit Runge--Kutta and spectral
deferred-correction methods.

\subsection{Statement}\label{statement}

Let \(s\ge2\) be an integer and \(0<c_1<\cdots<c_s\le1\) be real
collocation nodes. Let

\[\ell_j(t)=\prod_{k\ne j}\frac{t-c_k}{c_j-c_k},
\qquad A_{ij}=\int_0^{c_i}\ell_j(t)\,dt,
\qquad A\in\mathbb R^{s\times s}.\]

\textbf{Conjecture (van der Houwen--de Swart).} For every such node set,
does there exist a real positive diagonal matrix
\(D=\operatorname{diag}(d_1,\ldots,d_s)\) such that

\[\rho(I_s-D^{-1}A)=0?\]

Equivalently, require

\[ (I_s-D^{-1}A)^s=0,\]

or the polynomial identity

\[\det((1-t)D+tA)=\det D\qquad\text{for every }t\in\mathbb R.\]

No ordering constraint on the positive \(d_i\) is imposed. The source
states the conjecture for collocation-based Runge--Kutta correctors; the
explicit positive-distinct-node quantification here is an editorial
formalization of its collocation-matrix setting, also used by the 2025
treatment. The exclusion of a node at zero is essential: such a node
gives \(A\) a zero row and forces \(I_s-D^{-1}A\) to have eigenvalue
one. Rescaling positive nodes to \(c_s\le1\) does not change the
existence question.

\newpage
\subsection{Numerical significance and known
distinctions}\label{numerical-significance-and-known-distinctions}

For the diagonal-preconditioned iteration of a collocation system on the
scalar test equation, the iteration matrix is

\[K(z)=z(I_s-zD)^{-1}(A-D).\]

Its limit as \(|z|\to\infty\) is \(I_s-D^{-1}A\). The conjecture thus
asks whether the stiff-limit iteration can terminate in at most \(s\)
steps using a positive diagonal preconditioner, which permits
independent stage solves. It does not assert uniform contraction for
finite \(z\) or absence of transient nonnormal growth.

The case \(s=2\) has an elementary verification. Write \(a=c_1<b=c_2\)
and choose

\[d_1=\frac{a(2b-a)}{2(b-a)+\sqrt{2ab}},\qquad d_2=\frac{ab}{2d_1}.\]

Both entries are positive, and direct substitution gives
\(\operatorname{tr}(D^{-1}A)=2\) and \(\det(D^{-1}A)=1\).
Cayley--Hamilton therefore gives \((I_2-D^{-1}A)^2=0\). This
verification is included here explicitly; it is not a claim of an
all-stage result in the cited papers.

The 2025 MIN-SR-S construction searches for this nilpotence through
determinant equations, and explicitly leaves existence at arbitrary
stage count unproved. Its theorem on nilpotence of \(A-D\) concerns the
nonstiff limit, a different matrix. Its variable-sweep MIN-SR-FLEX
coefficients also do not supply one fixed \(D\) solving the displayed
target. Triangular preconditioners can achieve stiff-limit nilpotence by
LU splitting, but do not meet the diagonal constraint.

\subsection{References and status
check}\label{references-and-status-check}

\begin{itemize}
\tightlist
\item
  P. J. van der Houwen and J. J. B. de Swart, \emph{Triangularly
  implicit iteration methods for ODE-IVP solvers}, SIAM Journal on
  Scientific Computing 18(1) (1997), 41--55.
  \href{https://doi.org/10.1137/S1064827595287456}{DOI};
  \href{https://ir.cwi.nl/pub/2191/2191D.pdf}{CWI PDF}. §3.2.1, p.46,
  states the positive-diagonal zero-spectral-radius conjecture; §3.2.2,
  equation (3.9) and Theorem 3.1, specifies collocation matrices with
  positive distinct abscissas.
\item
  G. Čaklović, T. Lunet, S. Götschel and D. Ruprecht, \emph{Improving
  efficiency of parallel across the method spectral deferred
  corrections}, SIAM Journal on Scientific Computing 47(1) (2025),
  A430--A453. \href{https://doi.org/10.1137/24M1649800}{DOI and full
  text}; \href{https://d-nb.info/1363153935/34}{open PDF}. §2.2,
  equations (2.11)--(2.13), distinguishes the two limits; §2.2.3,
  p.A439, Definition 2.10 and equation (2.38), explicitly discusses
  unproved existence for arbitrary stage counts.
\end{itemize}

On 2026-09-11, checked both full papers, searched the original title and
positive diagonal/nilpotent collocation conjecture, and searched
MIN-SR-S existence, proofs and 2026 follow-ups. No theorem or
counterexample settling this fixed-diagonal existence target was
located. The 2025 paper gives current explicit evidence that the issue
remains unresolved; its numerical optimization is not an all-node,
all-stage existence proof. This is a bounded check. The target is
distinct from the
\href{https://github.com/ajt60gaibb/OpenProblemsInNLA/blob/main/linear-systems-and-elimination/IE-27/README.md}{LU-based
spectral disk question} and from existing incomplete-factorization
entries.
\end{document}
